\changecolours{ScoutGreen}
\section[Conclusão e Próximos Passos]{Conclusão, Atividades e Conexão com a Aula 07}

% ------------------------------------------------------------------------------
\twocolframe[title={Síntese Conceitual da Aula 06},leftcol=7.1cm,rightcol=6.9cm]{%
  \alert{Checklist de Competências Consolidadas}\par
  \itemseps[topsepi=0.5mm,itemsepi=1mm,labelsep=2mm,leftmargini=3mm]
  \begin{itemize}
    \item \textbf{Entidade Fraca:} Ausência de chave própria e dependência existencial/identificação.
    \item \textbf{Discriminador:} Chave parcial grafada com sublinhado tracejado (\dashuline{atributo}).
    \item \textbf{Gramática Chen:} Retângulo duplo, losango duplo e linha dupla para a entidade fraca.
    \item \textbf{Modelo EER:} Superclasses, subclasses e herança universal de atributos e relacionamentos.
    \item \textbf{Restrições EER:} Disjunção ($d/o$) e Integralidade (Total/Parcial).
  \end{itemize}
}{%
  \alert{Boas Práticas de Modelagem Conceitual}\par
  \itemseps[topsepi=0.5mm,itemsepi=1mm,labelsep=2mm,leftmargini=3mm]
  \begin{itemize}
    \item \textbf{Não Abusar de Entidades Fracas:}
      Só utilize entidade fraca se a existência e identificação forem estritamente dependentes de outra entidade do minimundo.
    \item \textbf{Clareza na Especialização:}
      Não crie subclasses se não houver atributos específicos ou relacionamentos exclusivos justificando a divisão.
    \item \textbf{Documentação Formal:}
      Sempre explicite se a hierarquia é ($d$) ou ($o$), e se a participação na superclasse é total ou parcial.
  \end{itemize}
}

% ------------------------------------------------------------------------------
\twocolframe[title={Ponte Pedagógica: Aula 07 (Sábado Letivo)},leftcol=7.1cm,rightcol=6.9cm]{%
  \alert{O Próximo Encontro: Aula 07 (17/10/2026)}\par
  \itemseps[topsepi=0.5mm,itemsepi=1mm,labelsep=2mm,leftmargini=3mm]
  \begin{itemize}
    \item \textbf{Sábado Letivo -- Estudo Dirigido AVP1:}
      Desenvolvimento do projeto prático de modelagem conceitual de banco de dados.
    \item \textbf{Escopo do Projeto Semestral:}
      Definição do minimundo de aplicação (gestão de frotas, reservas hoteleiras ou controle comercial).
    \item \textbf{Modelagem Hands-on:}
      Construção completa do DER conceitual aplicando entidades regulares, fracas, atributos e relacionamentos.
  \end{itemize}
}{%
  \alert{Atividades de Fixação da Aula 06}\par
  \itemseps[topsepi=0.5mm,itemsepi=1mm,labelsep=2mm,leftmargini=3mm]
  \begin{itemize}
    \item \textbf{Exercícios Conceituais:}
      Resolver as 5 questões analíticas do plano de aula no caderno de estudos.
    \item \textbf{Jupyter Notebook Interativo:}
      Executar o script Python para testar os bloqueios de discriminador duplicado e disjunção.
    \item \textbf{Esboço da AVP1:}
      Identificar ao menos uma entidade fraca ou hierarquia EER no escopo da sua equipe.
  \end{itemize}
}
